\documentclass[11pt]{article}

\usepackage[margin=1in]{geometry}
\usepackage{amsmath,amssymb}
\usepackage{bm}
\usepackage{booktabs}
\usepackage{enumitem}
\usepackage{hyperref}

\title{Differential Splicing with Subisoform Covariance}
\author{}
\date{}

\begin{document}
\maketitle

\section*{Purpose}

The hierarchical model estimates transcript abundance from equivalence-class
(EC) counts, but full-length transcripts are often less identifiable than the
local splicing choices needed for biological interpretation. The implemented
differential-splicing layer therefore collapses transcripts to local splice
paths, estimates covariance from the same paired-primer EC likelihood, and
tests path usage across biological replicates.

The analysis distinguishes two sources of variation:
\begin{enumerate}[leftmargin=*]
  \item inferential uncertainty caused by finite UMIs and ambiguous ECs;
  \item biological variation among mice with the same condition and cell type.
\end{enumerate}
Inferential covariance is not used as a replacement for biological
replication. It determines the precision assigned to each pseudobulk, while a
separate variance component represents between-mouse variation.

\section*{Subisoforms}

\subsection*{Atomic exon annotation}

Annotated exon coordinates are read from the GTF as zero-based, half-open
intervals. Overlapping exons within each transcript are merged. All exon
boundaries in a gene define atomic intervals. An atomic interval is
\emph{constitutive} when it is exonic in every annotated spliced transcript
of that gene. Adjacent constitutive intervals are merged into constitutive
anchor regions.

\subsection*{Splice-path blocks}

Each pair of consecutive constitutive anchors defines one block. A
transcript's subisoform in that block is its ordered set of exonic intervals
between the anchors. Transcripts with the same local path are collapsed even
if their full-length structures differ elsewhere. Blocks with only one path
are discarded.

Terminal pseudo-anchors are added before the first and after the last
constitutive region. A gene with no shared constitutive region is represented
by one whole-gene block. These rules include alternative first and last
regions while retaining the local interpretation of internal blocks. Simple
two-path blocks correspond to familiar skipped-exon, mutually-exclusive-exon,
alternative donor, alternative acceptor, and retained-intron events. The
general representation also permits more than two paths and nested events.

Let \(b\) index blocks in gene \(g\), and let \(C_{gb}\) be its
path-by-transcript incidence matrix. Precursor features are retained in the EC
likelihood as nuisance features but are not assigned to spliced paths.

\section*{Paired EC Likelihood}

For one mouse-by-cell-type pseudobulk and primer \(p\), let
\(\bm c_p\) be the EC counts, \(N_p=\bm 1^\prime\bm c_p\), and
\(A_p\) the fixed weighted EC design. The transcript abundance vector
\(\bm\theta\) gives
\[
  \bm q_p=A_p\bm\theta,\qquad
  Z_p=\bm 1^\prime\bm q_p=\bm d_p^\prime\bm\theta,
  \qquad
  \bm d_p=A_p^\prime\bm 1.
\]
Conditioning on \(N_p\),
\[
  \bm c_p\mid N_p,\bm\theta
  \sim
  \operatorname{Multinomial}
  \left(N_p,\frac{\bm q_p}{Z_p}\right).
\]
The oligo(dT) and random-hexamer likelihoods are evaluated separately and
their information matrices are added. The weighted designs preserve their
different positional and transcript-exposure assumptions.

Only ECs whose compatible transcripts map uniquely to one gene enter that
gene's covariance calculation. Cross-gene ECs remain part of the genome-wide
hierarchical fit but are excluded from this local curvature approximation.

\section*{Transcript Information}

Write \(\eta_t=\log\theta_t\). Define the EC-specific transcript
responsibilities and the primer-level sampled-transcript distribution
\[
  \bm r_{ep}
  =
  \frac{A_{ep\cdot}\odot\bm\theta}{q_{ep}},
  \qquad
  \bm r_{0p}
  =
  \frac{\bm d_p\odot\bm\theta}{Z_p}.
\]
The expected Fisher information in transcript-log-abundance coordinates is
\[
  I_{\eta,p}
  =
  N_p\left[
    \frac{
      \operatorname{diag}(\bm\theta)
      A_p^\prime\operatorname{diag}(\bm q_p^{-1})
      A_p\operatorname{diag}(\bm\theta)
    }{Z_p}
    -
    \bm r_{0p}\bm r_{0p}^\prime
  \right].
\]
Equivalently, this is \(N_p\) times the variance across observed ECs of the
conditional transcript responsibility. Ambiguous ECs have similar
responsibilities under competing transcripts and therefore contribute little
information in those contrast directions. The total information is
\[
  I_\eta=\sum_p I_{\eta,p}.
\]
This is the expected-information counterpart of the missing-information
identity used to obtain uncertainty after EM.

\section*{Path Coordinates and Covariance}

For block \(b\), spliced transcript abundances are normalized within their
gene and collapsed:
\[
  \bm\psi_b
  =
  C_b
  \frac{\bm\theta_{\mathrm{spliced}}}
       {\bm 1^\prime\bm\theta_{\mathrm{spliced}}}.
\]
An orthonormal Helmert basis \(Q_b\), satisfying
\(Q_b^\prime\bm 1=0\), defines identifiable log-ratio coordinates
\[
  \bm y_b=Q_b^\prime\log\bm\psi_b.
\]
If \(J_b=\partial\bm y_b/\partial\bm\eta^\prime\), the profile-style
delta covariance is
\[
  V_b
  \approx
  J_b I_\eta^+ J_b^\prime,
\]
where \(I_\eta^+\) is the inverse on the positive-eigenvalue subspace.
The implementation explicitly projects \(J_b\) onto the null space of
\(I_\eta\). A block is reported as unidentified rather than assigning zero
variance when its path contrast overlaps an information-null direction.

This calculation profiles all transcript directions, including differences
among full-length transcripts assigned to the same local path. It can
therefore yield finite path covariance even when individual transcripts are
not identifiable.

\section*{Local Path Refinement}

The hierarchical cell decoder supplies a baseline for each pseudobulk. Cell
transcript proportions are weighted by their UMI totals and aggregated within
mouse, cell type, and condition. For a block, local parameters
\(\bm\delta_b\) multiply all transcripts carrying the same path:
\[
  \theta_t(\bm\delta_b)
  \propto
  \theta_t^{(0)}
  \exp\{Q_{b,s(t)}^\prime\bm\delta_b\}.
\]
The total spliced mass and the transcript mixture within each path remain
fixed. L-BFGS optimizes the exact paired multinomial likelihood. Under this
parameterization,
\[
  \bm y_b(\bm\delta_b)=\bm y_b^{(0)}+\bm\delta_b.
\]
The inverse Fisher information for \(\bm\delta_b\) gives a second,
\emph{conditional} covariance. It is useful for optimization diagnostics and
parametric bootstrap validation but is less conservative than the profile
covariance because within-path transcript mixtures are fixed.

Finite covariance alone is not sufficient for a Gaussian approximation.
Weakly informed path fits can move to a boundary under repeated sampling.
The implemented reliability rule requires
\[
  \lambda_{\max}(V_{bi})\leq 1/8
  \quad\text{and}\quad
  \min_s\widehat\psi_{bis}\geq 0.15,
\]
so every retained orthonormal log-ratio direction has effective Fisher
information of at least eight and no fitted path lies near the simplex
boundary. All fits remain in the estimate output, with separate identifiable
and reliable flags; only reliable fits enter the differential test. The
thresholds were selected from the label-free parametric bootstrap described
below.

\section*{Differential-Splicing Test}

Let \(i\) index independent mouse pseudobulks within a cell type. For each
block,
\[
  \widehat{\bm y}_{bi}
  =
  X_i B_b+\bm\epsilon_{bi},
  \qquad
  \operatorname{Var}(\bm\epsilon_{bi})
  =
  V_{bi}+\tau_b^2 I.
\]
Here \(V_{bi}\) is the EC-derived inferential covariance. The scalar
\(\tau_b^2\) represents biological variation not explained by condition.
It is estimated by restricted maximum likelihood under the null design,
which excludes the condition coefficients. Generalized least squares combines
pseudobulks using
\((V_{bi}+\tau_b^2I)^{-1}\). A multivariate Wald test evaluates all
non-reference condition coefficients. Benjamini--Hochberg correction is
performed separately for conditional and profile covariance results.

Condition labels are permuted while preserving their observed group sizes.
The null REML variance is invariant to these permutations and is reused.
Permutation p-values and the distribution of permuted asymptotic p-values
are written as calibration diagnostics. The primary p-value is the
asymptotic Wald value because a small fixed number of permutations cannot
resolve the tail needed for genome-wide false-discovery control.

For a two-path block, the test has one log-ratio coordinate and is equivalent
to a precision-weighted differential-PSI analysis. Multi-path blocks receive
one omnibus test with \(S_b-1\) coordinates per condition contrast.

\section*{Validation}

Numerical tests cover:
\begin{itemize}[leftmargin=*]
  \item construction of a skipped-exon block from a synthetic GTF;
  \item finite path covariance when two indistinguishable transcripts collapse
        to one path;
  \item rejection of a path contrast when all paths produce the same EC;
  \item equality of the analytic two-path Fisher variance and the binomial
        result;
  \item agreement between analytic variance and parametric multinomial
        simulations;
  \item recovery of an imposed multivariate condition effect by GLS.
\end{itemize}

The production validation additionally simulates EC counts from fitted
pseudobulk path models, refits the local path parameters, and compares
empirical squared error with analytic conditional variance. It reports
component-wise squared standardized error and 95\% interval coverage for
each validation fit. Condition-label permutations provide a null calibration
check for the complete differential-testing layer.

In a two-path smoke analysis, all 1,257 eligible pseudobulk fits converged.
Conditional covariance was identifiable for 95.5\% of fits and profile
covariance for 52.3\%. Before reliability filtering, fits with information
below eight and paths below 15\% usage showed boundary-driven mode switching.
Above both thresholds, mean squared standardized error was 1.09 and
component-wise 95\% coverage was 94.0\%. Independent normal simulations of
the REML--GLS implementation gave
type-I error between 2.7\% and 5.0\% at a nominal 5\% across zero, moderate,
and large biological-variance settings.

The genome-wide microglialess analysis considered blocks with up to eight
paths. After the gene-count filter, 2,820 blocks contributed 46,393
pseudobulk estimates; every local optimization converged. Conditional
covariance was identifiable for 73.5\% of estimates and 12.0\% passed both
reliability criteria. Profile covariance was identifiable for 29.8\% and
7.9\% passed both criteria. The retained estimates produced 170 conditional
and 113 profile differential tests. Five conditional and three profile tests
had unadjusted \(p<0.05\), but none passed 5\% false-discovery control.

Across 3,000 production bootstrap refits, mean squared standardized error was
1.52 and 95\% interval coverage was 94.6\%. Among permuted-label tests, 2.9\%
of conditional and 2.5\% of profile asymptotic p-values were below 0.05.
Thus the implemented filter gives slightly conservative null calibration.
The lack of discoveries reflects both this conservatism and the small number
of mouse replicates available for many cell-type--condition combinations.

\section*{Limitations}

The profile covariance is conditional on the fitted genome-wide decoder.
Uncertainty in shared decoder parameters is expected to be small for common
paths but may matter for rare genes or cell types. A mouse-level jackknife of
the hierarchical fit is the preferred extension.

The current biological variance is isotropic in log-ratio space. Estimating a
full covariance separately for every block would be unstable with few mice.
A later empirical-Bayes model could shrink path-specific biological
covariances across genes.

\end{document}
